\subsection{\texorpdfstring{关系式 \(\sum_{i} e_i f_i \leqslant n\)}{关系式 Σ ei fi ≤ n}}

设 K 为一个域，v 为 K 上的一个赋值，L 为 K 的一个次数为 n 的有限扩张。~设 \((v_{1}^{\prime}, \ldots, v_{r}^{\prime})\) 为 L 上延拓 v 的赋值，其中任意两个均不等价；若它们独立（当 v 的高为 1 时总是如此），则

\[
\sum_ {i = 1} ^ {r} e (v _ {i} ^ {\prime} / v) f (v _ {i} ^ {\prime} / v) \leqslant n
\]

由命题 2（公式 (6) 和 (7)）可知。我们将看到，这个结果在一般情形下仍然成立。准确地说：

定理 1. 设 K 为一个域，v 为 K 上的一个赋值，L 为 K 的一个次数为 n 的有限扩张。~则：

\begin{enumerate}
\def\labelenumi{(\alph{enumi})}
\item
  每个完备延拓系 \((v_i')_{i\in I}\) 的 \(v\) 到 \(\mathbf{L}\) 的延拓都是有限的。
\item
  \(\sum_{i\in I}e(v_i^{\prime}|v)f(v_i^{\prime}|v)\leqslant n\) 从而 Card(I) \(\leqslant n.\)
\item
  任意两个 \(v_{i}'\) 的环关于包含关系均不可比较。
\end{enumerate}

由于 \(v\) 为非真赋值时定理显然成立，我们设 \(v\) 不为非真赋值。设 \((v_1', \ldots, v_s')\) 为 \(L\) 上延拓 \(v\) 的任意有限族，其中任意两个均不等价。我们首先证明 \(\sum_{i=1}^{s} e(v_i'/v)f(v_i'/v) \leqslant n\)。这将证明 (a) 和 (b)。

我们对 s 归纳，并因此假设不等式已对 0、1、\ldots、s - 1 个赋值的情形成立。分两种情形。

\begin{enumerate}
\def\labelenumi{(\arabic{enumi})}
\tightlist
\item
  假设至少存在两个独立赋值 \(v_i'\)。则存在（§ 7，第 2 节，注 1）一个 \([1, s] = I_1 \cup \cdots \cup I_t\) 对 \([1, s]\) 的划分，使得：
\end{enumerate}

\begin{enumerate}
\def\labelenumi{(\roman{enumi})}
\item
  赋值 \(v_{i}^{\prime}\) 和 \(v_{j}^{\prime}\) 相关的充要条件是 \(i\) 和 \(j\) 属于同一个 \(I_{k}\)；
\item
  \(\operatorname{Card}(\mathbf{I}_k) < s\) 对所有 \(k\) 均成立。
\end{enumerate}

我们在每个 \(I_{k}\) 中选取一个指标 \(i(k)\)。令 \(\hat{L}_{i(k)}\) 表示 L 关于 \(v_{i(k)}^{\prime}\) 的完备化，并令 \(n(k) = [\hat{L}_{i(k)} : \hat{K}]\)。对于所有 \(i \in I_{k}\)，\(v_{i}^{\prime}\) 在 L 上定义出与 \(v_{i(k)}^{\prime}\) 相同的拓扑（§ 7，第 2 节，命题 3），因而可将其延拓为赋值 \(\hat{v}_{i}^{\prime}\)，定义在 \(\hat{L}_{i(k)}\) 上，其在 \(\hat{K}\) 上的限制是 \(\hat{v}\)。由于 \(v_{i}^{\prime}\)（\(i \in I_{k}\)）中任意两个均不等价，\(\hat{v}_{i}^{\prime}\) 也具有同样性质。将归纳假设应用于有序对（\(\hat{K}, \hat{L}_{i(k)}\)），根据命题 2 (a) 的公式 (4)，得到 \(\sum_{i\in I_k}e(v_i^{\prime}|v)f(v_i^{\prime}|v)\leqslant n(k)\)。由于 \(\sum_{k = 1}^{t}n(k)\leqslant n\)（命题 2 (b)，公式

(7)），当然有 \(\sum_{i=1}^{s} e(v_i'/v) f(v_i'/v) \leqslant n\)。

\begin{enumerate}
\def\labelenumi{(\arabic{enumi})}
\setcounter{enumi}{1}
\tightlist
\item
  现在转到任意两个 \(v_i'\) 都相关的情形。令 \(A_i'\) 为 \(v_i' (1 \leqslant i \leqslant s)\) 的环；记 A 为赋值 \(v, A_i' \cap K = A\) 的环，对所有 \(i\) 都成立。令 \(B'\) 为由 \(A_1', \ldots, A_s'\) 生成的 L 的子环；记 \(B = B' \cap K\)；则 \(B \supset A\)。于是 B 是 K 上赋值 \(w\) 的环，而 \(B'\) 是赋值 \(w'\) 的环，该赋值非真且延拓 \(w\)（\(\S 7\)，第 2 节，命题 4）；域 \(\kappa(B')\) 是 \(\kappa(B)\) 的次数为 \(f(w'/w)\) 的扩张。考察典范像 \(\bar{A}_i', \bar{A}\)，它们是 \(A_i'\) 和 A 在 \(\kappa(B')\) 中的像；则 \(\bar{A}\) 是赋值 \(\bar{v}\) 在 \(\kappa(B)\) 上的环，而 \(\bar{A}_i'\) 是 \(\bar{v}_i'\) 在 \(\kappa(B')\) 上延拓 \(\bar{v}\) 的环。由于 \(A_i'\) 生成 \(B'\)，\(\bar{A}_i'\) 生成 \(\kappa(B')\)，所以 \(\bar{v}_i'\) 不全相关（\(\S 7\)，第 2 节，命题 4）。由证明的第一部分，
\end{enumerate}

\[
\sum_ {i = 1} ^ {s} e \left(\bar {v} _ {i} ^ {\prime} / \bar {v}\right) f \left(\bar {v} _ {i} ^ {\prime} / \bar {v}\right) \leqslant \left[ \kappa \left(\mathbf {B} ^ {\prime}\right): \kappa (\mathbf {B}) \right] = f \left(w ^ {\prime} / w\right)
\]

从而

\[
\sum_ {i = 1} ^ {s} e (w ^ {\prime} / w) e (\bar {v} _ {i} ^ {\prime} / \bar {v}) f (\bar {v} _ {i} ^ {\prime} / \bar {v}) \leqslant e (w ^ {\prime} / w) f (w ^ {\prime} / w) \leqslant n \quad \text {(no. 1, Lemma 1)}.
\]

因此，只要证明

\[
f (\bar {v} _ {i} ^ {\prime} / \bar {v}) = f (\bar {v} _ {i} ^ {\prime} / \bar {v}), \quad e (w ^ {\prime} / w) e (\bar {v} _ {i} ^ {\prime} / \bar {v}) = e (v _ {i} ^ {\prime} / v).\tag{8}
\]

为此，注意到 v 和 \(\bar{v}\)（分别地，\(v_{i}^{\prime}\) 和 \(\bar{v}_{i}^{\prime}\)）具有相同的剩余域（§ 4，第 1 节，命题 2 的推论）；这证明了第一个等式。对于第二个等式，根据 § 4，第 3 节的注，有如下交换图，其中各行是正合列，竖直箭头表示典范嵌入：

\[
\begin{array}{c} 0 \to \Gamma_ {\bar {v}} \to \Gamma_ {v} \to \Gamma_ {w} \to 0 \\ \downarrow \quad \downarrow \quad \downarrow \\ 0 \to \Gamma_ {\bar {v} _ {i} ^ {\prime}} \to \Gamma_ {v _ {i} ^ {\prime}} \to \Gamma_ {w ^ {\prime}} \to 0 \end{array}
\]

由此得到正合列

\[
0 \rightarrow \Gamma_ {\bar {v} _ {i} ^ {\prime}} / \Gamma_ {\bar {v}} \rightarrow \Gamma_ {v _ {i} ^ {\prime}} / \Gamma_ {v} \rightarrow \Gamma_ {w ^ {\prime}} / \Gamma_ {w} \rightarrow 0
\]

由第一章第 1 节第 4 节命题 2，得到一个正合列，这证明了第二个公式 (8)。为了完成定理 1 的证明，还需证明 (c)。若 \(v_i'\) 的环包含 \(v_j'\) 的环，则 \(\Gamma_{v_i'}\) 可识别为商群 \(\Gamma_{v_j'} / H\)，其中 \(H\) 是孤立子群（§ 4，第 3 节）。由于复合的典范映射

\[
\Gamma_ {v} \rightarrow \Gamma_ {v _ {j} ^ {\prime}} \rightarrow \Gamma_ {v _ {j} ^ {\prime}} / H = \Gamma_ {v _ {i} ^ {\prime}}
\]

是单射，\(H \cap \Gamma_{v} = \{0\}\)，从而 \(H = \{0\}\)（引理 3，第 1 节）。于是 \(v_{i}'\) 和 \(v_{j}'\) 等价，故 i = j。

注。环 \(A_{i}^{\prime}\)（对应赋值 \(v_{i}^{\prime}\)，其中 \(i \in I\)）的交 C 是 A 在 L 中的整闭包（\(\S 1\)，第 3 节，定理 3 的推论 3）；此外由 (c) 以及 \(\S 7\)，第 1 节，命题 1 和 2 可知，C 是半局部环，其极大理想是交 \(m_{i} = C \cap m(A_{i}^{\prime})\)，并且 \(A_{i}^{\prime} = C_{m}\) 对所有 \(i \in I\) 都成立。

